Fix an integer $n\ge 2$.

Definition (admissible interval factorisation).
For integers $t\ge 2$ and $m\ge 1$, we say that $n!$ has an admissible $(t,m)$-factorisation if there exist integers
\[
1\le a_1 < a_2 < \cdots < a_t
\]
with
\[
a_t=a_1+m\quad\text{and}\quad n! = \prod_{i=1}^t a_i.
\]
Equivalently, there exist integers $A\ge 1$ and $m\ge 1$ and a subset $S\subseteq\{A,A+1,\dots,A+m\}$ containing both endpoints $A$ and $A+m$ such that
\[
\prod_{s\in S} s = n!.
\]

Definition (the function $f$).
For $n\ge 2$ define
\[
f(n):=\min\{m\ge 1:\; n!\text{ has an admissible }(t,m)\text{-factorisation for some }t\ge 2\}.
\]

The problem asks for the asymptotic behaviour of $f(n)$ as $n\to\infty$, and in particular whether $f(n)=1$ for infinitely many $n$.

For $m\ge 1$ and $N\ge 2$, define
\[
F_m(N):=|\{n\le N:\ f(n)=m\}|.
\]

Edge cases.
For $n=1$, the definition as stated has no solution because $1!=1$ cannot be written as a product of at least two distinct positive integers; throughout I treat $f$ as defined on $n\ge 2$.
